\section{The Cross-Cultural Variation of Flourishing and Its Practice}
\label{p11:sec:crosscultural}

\noindent
The account developed so far has proceeded at a level of generality that abstracts from the cultural particularity of the shared lives it concerns, and this abstraction, useful for the statement of structure, conceals a variation that a foundational treatment must address. The conception of flourishing, the practices through which a shared life is cultivated, the standards by which it is evaluated, and the very methods by which it is studied vary across cultures, and the variation is not noise around a universal mean but a systematic difference in what flourishing is taken to be and how it is pursued. This section addresses three strands of this variation: the cross-cultural variation in the conception of happiness and flourishing, the variation across scholarly traditions in the method by which flourishing is studied, and the variation across cultures in the practices of marriage and family. It then locates the variation within the lifecycle, treating each moment of the cycle as the site of a culturally specific cognition and practice. Before the strands are taken up, the basis on which cultures are selected and compared is stated, so that the comparison proceeds by a method rather than by an unexamined choice of cases. The section draws throughout on the operational concession and on the framework of symmetry-broken phases (\xref{p11:sec:bg-phases}), of which the cultural variation is a further and consequential instance.

\subsection{The basis of the comparison: method and case selection}
\label{p11:sec:cc-method-basis}

A comparison across cultures requires a principled basis for the selection of the cultures compared and for the dimensions along which they are compared, on pain of mistaking an arbitrary choice of cases for a representative sample of human variation. The comparative method in the social sciences supplies two complementary logics of case selection, both descending from Mill's methods of agreement and difference \parencite{mill1843} and systematised for social inquiry by \textcite{przeworski1970}. The most-similar-systems design selects cases alike in most respects and differing in the variable of interest, so that a difference in outcome may be attributed to that variable; the most-different-systems design selects cases differing in most respects, so that a shared outcome may be attributed to the few factors they hold in common \parencite{lijphart1971}. The two designs are the principal means of managing the comparison of a small number of complex cases.

The cultures to be compared are located on an empirically derived map of cultural variation, so that their selection may be assessed for coverage and representativeness rather than asserted. The most developed such map is the one derived by Inglehart and Welzel from the World Values Survey, which resolves cross-national variation into two principal dimensions, a dimension from traditional to secular-rational values and a dimension from survival to self-expression values, and which locates national cultures in clusters shaped by their religious and historical heritage, among them the Protestant European, the Catholic European, the English-speaking, the Confucian or East Asian, the Latin American, the Orthodox, the South Asian, and the African-Islamic \parencite{inglehart2005}. Comparable coordinate systems are supplied by Hofstede's dimensions of national culture, of which the individualism-collectivism dimension is the one most pertinent here \parencite{hofstede2001}, and by Schwartz's cultural value orientations \parencite{schwartz2006}. These maps make the selection of cases assessable: a selection covers the variation well to the extent that it spans the principal clusters, and a case represents its cluster to the extent that it is typical of it.

The present treatment selects its cases on this basis, and states the selection openly. Four cultures are taken as the cases of detailed analysis, China, Japan, Continental Europe, and the United States, and the selection is a combination of the two comparative designs. China and Japan are two cases within the East Asian, broadly Confucian cluster, alike in their interdependent heritage and differing in its present configuration, and Continental Europe and the United States are two cases within the Western cluster, alike in their independent heritage and differing in their institutional and methodological traditions; the within-cluster pairs constitute most-similar comparisons, and the between-cluster contrast of the Western with the East Asian pair constitutes a most-different comparison. The selection is, deliberately, weighted toward the cultures the present writer is best placed to treat with care, and it is for that reason a selection of cases for detailed analysis and not a claim to global coverage. Its limitation is therefore stated plainly: the four cases span the principal axis of individualism and interdependence but do not span the full cultural map, and several major clusters, the South Asian, the African-Islamic, the Latin American, and the Orthodox, enter the treatment only as points of contrast and not as cases of detailed analysis. Where they bear on the argument they are noted, in particular the South Asian institution of arranged marriage and the joint family, the Latin American combination of collectivism with an actively expressed positive emotionality that differs from the East Asian emotional balance, and the Islamic configuration of kinship obligation; but the detailed analysis remains that of the four core cases, and the conclusions are offered with the representativeness this selection affords and no more.

\subsection{The cross-cultural variation in the conception of flourishing}
\label{p11:sec:cc-happiness}

The conception of happiness itself varies across cultures, and the principal axis of variation is the one between an independent and an interdependent construal of the self and its flourishing. In the cultural-psychological tradition, the European-American model of the self is independent, defined in terms of individual attributes, autonomy, and the pursuit of personal goals, while the East Asian model is interdependent, defined in terms of social relationships and the adjustment of the self to a relational whole \parencite{markus1991}. Happiness is construed accordingly. In North American contexts, happiness tends to be defined in terms of personal achievement, and individuals are motivated to maximise positive affect and to seek happiness through autonomous agency; in East Asian contexts, happiness tends to be defined in terms of interpersonal connectedness and the balance of the self with others \parencite{uchida2004}. The difference extends to the structure of emotional experience: where the North American construal seeks to maximise positive over negative affect, the East Asian construal tends toward a dialectical balance in which positive and negative affect are held together rather than the former maximised, and in which happiness is understood as fraught with potential social consequence rather than as an unalloyed good to be pursued \parencite{uchida2009}. The interdependent construal has been developed into a distinct conception of interdependent happiness, predicated on interpersonal harmony, quiescence, and ordinariness rather than on personal high arousal \parencite{hitokoto2015}.

This variation bears directly on the conception of flourishing the present series has developed, and the bearing is instructive rather than merely cautionary. The series has construed flourishing as the unfolding of each party's own dynamics, sustained within a shared life whose value circulates and returns; the conception is neither purely independent, since the unfolding is sustained within and returns to a relation, nor purely interdependent, since it is the unfolding of each that is at issue and not the submersion of each in a relational whole. The conception occupies, deliberately, a position between the independent and interdependent construals, and the cross-cultural evidence locates it: it is a conception for which the relation is the medium of each party's development and the development is genuinely of each, which is neither the autonomous self-development of the independent model nor the relational adjustment of the interdependent one, but a structure in which the two are made compatible by the requirement that the relation return to each what each gives. The cultural variation is, in the terms of \xref{p11:sec:bg-phases}, a set of symmetry-broken phases of the cognition of flourishing, the independent and interdependent construals being two such phases, and the series' conception an attempt to articulate a structure of which both are partial determinations.

\subsection{The comparative variation in the method of study}
\label{p11:sec:cc-method}

A second variation, less often remarked in connection with the first but consequential for a foundational treatment, is the variation across scholarly traditions in the method by which flourishing and the shared life are studied. The methodological division between an explanatory, quantitative, positivist social science and an interpretive, qualitative, hermeneutic one is itself culturally and geographically patterned, and the pattern bears on the framework-divide of \xref{p11:sec:bg-phases}.

The interpretive tradition originates in the nineteenth-century German distinction between the natural and the human sciences, the \emph{Naturwissenschaften} and the \emph{Geisteswissenschaften}, and in the corresponding distinction between explanation, \emph{Erklären}, the method of the natural sciences, and understanding, \emph{Verstehen}, the method of the human sciences, drawn by Droysen and Dilthey and developed by Weber \parencite{dilthey1883,weber1904}. On this tradition, the social world is constituted by the meanings its participants give it, and its study requires the interpretive recovery of those meanings, the production of what Geertz called thick description, which penetrates the webs of significance within which human action is intelligible rather than reducing action to observable behaviour \parencite{geertz1973}. The positivist tradition, by contrast, holds that the social world is continuous with the natural, that it is to be studied by the observation and measurement of behaviour and the discovery of regularities, and that what cannot be brought to measurement cannot be known. The two traditions are unevenly distributed: the interpretive and hermeneutic tradition has been more central to Continental European social thought, while the positivist and quantitative tradition has been more dominant in American social science, a difference of scholarly culture that patterns how flourishing is studied on either side.

This methodological variation is, for the present paper, the framework-divide of \xref{p11:sec:bg-phases} appearing at the level of scholarly culture rather than of the individual investigator. The positivist framework and the phenomenological-hermeneutic framework, which the paper has treated as symmetry-broken phases of the study of flourishing, are not merely options available to any investigator but traditions institutionalised in different scholarly cultures, the one measuring what can be measured and the other interpreting what is meant. The recognition has a consequence for the paper's method. The polyphonic refusal to subordinate one framework to another (\xref{p11:sec:bg-phases}) is, read in this light, a refusal to subordinate one scholarly culture's method to another's, and the operational concession, which admits the structured measure while denying it the status of the real, is a way of drawing on the positivist tradition's instruments without adopting its claim that what those instruments measure exhausts what flourishing is. The paper's stance toward the methodological divide is the stance the cultural variation recommends: that the explanatory and the interpretive methods are partial determinations of the study of a flourishing that neither exhausts, and that a foundational treatment must draw on both without resolving their difference by the victory of one.

\subsection{The cross-cultural variation in the practices of marriage and family}
\label{p11:sec:cc-family}

A third variation concerns the practices of marriage and family through which a shared life is cultivated, and it is the variation most directly pertinent to a praxis of intimate flourishing. The conception of what a marriage is, what a family is, and what the parties owe one another varies across cultures along several axes, of which three are salient.

The first axis concerns the meaning and basis of marriage. A long sociological account traces, in the Euro-American context, a transition from the institutional marriage, founded on social obligation and the fulfilment of prescribed roles, through the companionate marriage of the mid-twentieth century, founded on the personal ties of affection and companionship, to the individualised marriage of the present, founded on personal choice and the self-development of the partners \parencite{cherlin2004}. The transition is held to have weakened the social norms that once defined the partners' behaviour, a process termed the deinstitutionalisation of marriage, though its extent varies by social class and the institution retains, especially among the more educated, a central place \parencite{cherlin2020}. The individualised marriage, in which the partners actively select one another, negotiate their commitment, and pursue self-development within the relation, is the form the European sociology of intimacy has described as the normal chaos of love and the transformation of intimacy \parencite{beck1995,giddens1992}.

The second axis concerns the relation of the couple to the wider family. In the more individualist Euro-American contexts, the married couple tends to be the primary unit, relatively detached from the extended family; in the more collectivist East and Southeast Asian contexts, the couple is embedded in an extended family to which it owes ongoing obligations, and the marriage is as much a relation between families as between individuals. The obligation of adult children to their parents, filial piety, is a central organising norm of the East Asian family, and recent work distinguishes a reciprocal filial piety, founded on affection and a lasting positive bond, from an authoritarian filial piety, founded on obedience to social obligation, the former proving more nearly universal across cultures and the latter more culturally specific \parencite{yeh2013}. The intergenerational solidarity of families, the structure of obligation and support across generations, is itself a subject of formal comparative study \parencite{bengtson1991}, and the autonomy-relatedness relation, on which the independent and interdependent construals divide, has been argued to admit a synthesis in which autonomy and relatedness are not opposed but jointly developed, a model of particular relevance to families in modernising societies \parencite{kagitcibasi2005}.

The third axis concerns the direction and pace of change. The family systems of East Asia have been described as at a crossroads, meeting Western individualism, modernity, and changing gender roles against a background of persisting traditional obligation, so that the contemporary East Asian family combines elements of the interdependent and the independent models in a configuration unlike either \parencite{ji2015}. The variation is therefore not static, a fixed map of cultural difference, but dynamic, a set of trajectories of change that differ in direction and pace across cultures. These variations are summarised in Table~\ref{p11:tab:family}.

\begin{table}[t]
\centering
\small
\renewcommand{\arraystretch}{1.25}
\caption{Axes of cross-cultural variation in the conception and practice of marriage and family.}
\label{p11:tab:family}
\begin{tabularx}{\linewidth}{@{}>{\raggedright\arraybackslash}p{2.7cm} >{\raggedright\arraybackslash}X >{\raggedright\arraybackslash}X >{\raggedright\arraybackslash}p{2.4cm}@{}}
\toprule
\rowcolor{pageblush}
\textbf{\textcolor{rosedeep}{Axis}} & \textbf{\textcolor{rosedeep}{Individualist pole (Euro-American)}} & \textbf{\textcolor{rosedeep}{Interdependent pole (East Asian)}} & \textbf{\textcolor{rosedeep}{Sources}} \\
\midrule
Basis of marriage & Individualised: personal choice, self-development; deinstitutionalised & Institutional and relational: obligation, family alliance, prescribed roles & \textcite{cherlin2004,cherlin2020} \\
Couple and wider family & Couple as primary unit, detached from extended family & Couple embedded in extended family with ongoing obligation & \textcite{ji2015,bengtson1991} \\
Intergenerational norm & Weak filial obligation; support voluntary & Filial piety, reciprocal and authoritarian; support obligatory & \textcite{yeh2013,kagitcibasi2005} \\
Conception of intimacy & Transformation of intimacy; the normal chaos of love & Harmony, ordinariness, interdependent happiness & \textcite{giddens1992,beck1995,hitokoto2015} \\
\bottomrule
\end{tabularx}
\end{table}

\subsection{The cultural variation across the moments of the cycle}
\label{p11:sec:cc-lifecycle}

The variations surveyed above are not external to the lifecycle but enter at each of its moments, and the foundational value of the account is increased by tracing how each moment is culturally inflected. The lifecycle furnishes a structure within which the cultural variation can be organised, and the organisation shows the variation to be systematic rather than diffuse.

At the moment of \emph{cognition}, the cultural variation is the variation in what flourishing is taken to consist in. The independent construal cognises the flourishing of a shared life as the mutual support of two self-developing individuals; the interdependent construal cognises it as the achievement of relational harmony and the fulfilment of relational obligation; and the conception of this series, the just unfolding of each within a returning cycle, cognises it as a structure of which both are partial. The cultural variation at this moment is a variation in the normative specification of \xref{p11:sec:cog-spec}, and a foundational treatment must hold that the specification, though it aims at a structure, is realised in culturally specific forms.

At the moment of \emph{practice}, the cultural variation is the variation in the habits, household forms, and family obligations through which a shared life is cultivated. The individualised marriage is practised through the negotiation of commitment and the mutual support of self-development; the interdependent family is practised through the fulfilment of filial and relational obligation and the maintenance of harmony; and the family manner cultivated, the \emph{habitus} of \xref{p11:sec:prac-virtue}, differs accordingly, the one forming dispositions of autonomous partnership, the other dispositions of relational embeddedness. The cultivation-not-legislation constraint of \xref{p11:sec:prac-constraint} takes a culturally specific form: what counts as the legislation of a habit, and what as its cultivation, depends on the cultural understanding of the family's authority, so that a practice that is cultivation in one cultural context may be experienced as legislation in another.

At the moment of \emph{evaluation}, the cultural variation is the variation in the standards by which a shared life is judged to be going well, and it bears on the measurement of \xref{p11:sec:evaluation}. The instruments surveyed there were developed largely within the independent, Euro-American construal, and they measure largely what that construal counts as flourishing, the individual's satisfaction and self-development; applied across cultures, they may systematically mismeasure the flourishing of interdependent shared lives, registering as deficient a relational harmony that the interdependent construal counts as flourishing, or registering as flourishing an autonomous self-development that the interdependent construal counts as a failure of relational obligation. The operational concession bears with particular force here: the structured measure is a proxy not only for an unsymbolised value but for a culturally specific construal of that value, and its cross-cultural application requires a caution beyond the one the concession already imposes.

At the moments of \emph{reflection} and \emph{re-cognition}, the cultural variation is the variation in how a shared life revises its sense of itself and transmits what it has learned. The transmission of a family manner across generations, the \emph{habitus} inherited and revised, is more central to the interdependent family, in which the marriage is embedded in a lineage, than to the individualised marriage, in which the couple is relatively detached; and the re-cognition that closes a turn of the cycle is, in the interdependent context, a re-cognition conducted with reference to the extended family and its continuity, where in the individualist context it is conducted more nearly between the partners alone. The holonomy of the cycle, the phase accumulated across its turns, is in the one case the development of a lineage and in the other the development of a partnership, and the difference is a difference in the unit across which the flourishing is reckoned. Table~\ref{p11:tab:cc-cycle} summarises the cultural inflection of the moments of the cycle.

\begin{table}[t]
\centering
\small
\renewcommand{\arraystretch}{1.25}
\caption{The cultural inflection of the moments of the lifecycle: how each moment differs under the independent and interdependent construals.}
\label{p11:tab:cc-cycle}
\begin{tabularx}{\linewidth}{@{}>{\raggedright\arraybackslash}p{2.4cm} >{\raggedright\arraybackslash}X >{\raggedright\arraybackslash}X@{}}
\toprule
\rowcolor{pageblush}
\textbf{\textcolor{rosedeep}{Moment}} & \textbf{\textcolor{rosedeep}{Independent construal}} & \textbf{\textcolor{rosedeep}{Interdependent construal}} \\
\midrule
Cognition & Flourishing as mutual support of two self-developing individuals & Flourishing as relational harmony and fulfilment of obligation \\
Practice & Negotiation of commitment; cultivation of autonomous partnership & Fulfilment of filial and relational obligation; embeddedness \\
Evaluation & Individual satisfaction and self-development; existing instruments fit & Relational harmony; existing instruments may misregister \\
Reflection & Conducted between the partners; standard is the partnership & Conducted with reference to the extended family and lineage \\
Re-cognition & Holonomy as development of a partnership & Holonomy as development of a lineage \\
\bottomrule
\end{tabularx}
\end{table}

\subsection{The standing of the cross-cultural account}
\label{p11:sec:cc-standing}

The cross-cultural variation does not refute the structural account but locates it, and the relation between the two requires statement. The account of the lifecycle, of the moments and their dynamics, of the cycle's persistence and phase, and of justice as the internal condition of the good cycle, is offered as a structure, and the cultural variation is the variation in the determinate forms the structure assumes. The independent and interdependent construals are two such forms, two symmetry-broken phases of the cognition and practice of a shared flourishing, and the series' conception, the just unfolding of each within a returning cycle, is offered not as a third culture's construal set beside them but as an attempt to articulate the structure of which the cultural construals are partial determinations. The attempt is itself culturally situated, written from a particular location and formed by particular traditions, and it claims neither a view from nowhere nor a neutrality among cultures, which would be the metalanguage the series denies. It claims only to articulate, from its situation, a structure that the cultural variation instances, and to hold its own articulation open to the correction that the encounter with each cultural form supplies. The justice condition, in particular, is offered across the cultural variation: that the value a party generates return to that party, and that no party be reduced to the sustaining fuel of a relation that returns to it nothing, is a condition a shared life may fail in any culture, by the culturally specific means each culture affords, and the skew of the return distribution (\xref{p11:sec:ev-skew}) is a proxy for an injustice that the independent and the interdependent family are each, in their own forms, capable of.
